\documentclass[10pt]{mypackage}

\usepackage[uprightgreeks,light]{kpfonts}
\renewcommand{\coloneq}{\coloneqq}
%\renewcommand{\mathbb}{\mathds}
\usepackage{homework}
%\usepackage{notes}

%\usepackage[ backend=bibtex, style = alphabetic, sorting=ynt ]{biblatex}
%\addbibresource{  }

\usepackage{parskip}

\fancyhf{}
\fancyhead[R]{Avinash Iyer}
\fancyhead[L]{Algebraic Topology: Homework 27}
\fancyfoot[C]{\thepage}

\setcounter{secnumdepth}{0}

\begin{document}
\RaggedRight
\section{Revised Problems}%
\begin{problem}[Homework 25, Problem 1 (b)]
  Prove that the inclusion $f\colon \left( D^n,S^{n-1} \right)\rightarrow \left( D^n,D^n\setminus\set{0} \right)$ is not a homotopy equivalence of pairs, in that there is no map $g$ such that $g\circ f$ and $f\circ g$ are homotopic to identity maps via maps of pairs.
\end{problem}
\begin{solution}
  Consider the annular subset of $\left( D^n,D^n\setminus\set{0} \right)$ given by
  \begin{align*}
    U &= \set{x\in \R^{n} | 1/2\leq \left\vert x \right\vert\leq 3/4}
  \end{align*}
  that is contained in $D^n\setminus\set{0}$ and has $U = \overline{U}\subseteq \left( D^n\setminus\set{0} \right)^{\circ}$. By excision, this induces an isomorphism
  \begin{align*}
    H_{\ast}\left( D^n,D^n\setminus\set{0} \right)\cong H_{\ast}\left( D^n\setminus U,\left( D^n\setminus \set{0} \right)\setminus U \right).
  \end{align*}
  Via a deformation retraction, we have that $\left( D^n\setminus \set{0} \right)\setminus U\simeq S^{n-1}\sqcup S^{n-1}$.

  From the long exact sequence on the pair $\left( D^n,S^{n-1} \right)$, we have
  \begin{align*}
    \widetilde{H}_{n}\left( D^n \right)\rightarrow \widetilde{H}_{n}\left( D^{n},S^{n-1} \right) \rightarrow \widetilde{H}_{n-1} \left( S^{n-1} \right) \rightarrow \widetilde{H}_{n-1}\left( D^{n} \right).
  \end{align*}
  Since $D^n$ is contractible, exactness gives that $\widetilde{H}_n\left( D^{n},S^{n-1} \right)\cong \widetilde{H}_{n-1}\left( S^{n-1} \right)$. Similarly, replacing $S^{n-1}$ with the homotopy equivalent $D^{n}\setminus D^{n}\setminus\set{0}$, we have that $ \widetilde{H}_n\left( D^n,D^n\setminus\set{0} \right)\cong \widetilde{H}_{n-1}\left( D^n\setminus\set{0} \right) $, and by excision, we have 
  \begin{align*}
    \widetilde{H}_n\left( D^n,D^n\setminus\set{0} \right) &\cong \widetilde{H}_n\left( D^n\setminus U,\left( D^n\setminus\set{0} \right)\setminus U \right).
  \end{align*}
  Finally, in the long exact sequence of $\left( D^n\setminus U,\left( D^n\setminus\set{0} \right)\setminus U \right)$, we have
  \begin{align*}
    \widetilde{H}_n\left( D^n\setminus U \right)\rightarrow \widetilde{H}_n\left( D^n\setminus U,\left( D^n\setminus\set{0} \right)\setminus U \right)\rightarrow \widetilde{H}_{n-1}\left( \left( D^n\setminus\set{0} \right)\setminus U \right)\rightarrow \widetilde{H}_{n-1}\left( D^n\setminus U \right).
  \end{align*}
  Since $D^n$ and $D^n\setminus U$ are respectively contractible and a disjoint union of contractible spaces, it follows that we have an isomorphism from 
  \begin{align*}
    \widetilde{H}_n\left( D^n\setminus U,\left( D^n\setminus\set{0} \right)\setminus U \right) &\cong \widetilde{H}_{n-1}\left( \left( D^n\setminus\set{0} \right)\setminus U \right).
  \end{align*}
  Suppose toward contradiction that a homotopy inverse for $f$ exists. This then induces the family of isomorphisms given by
  \begin{align*}
    \Z &\cong \widetilde{H}_{n-1}\left( S^{n-1} \right)\\
    &\cong \widetilde{H}_{n}\left( D^n,S^{n-1} \right)\\
                                              &\xrightarrow{f_{\ast}} \widetilde{H}_n\left( D^n,D^n\setminus\set{0} \right)\\
                                              &\rightarrow \widetilde{H}_n\left( D^n\setminus U,\left( D^n\setminus\set{0} \right)\setminus U \right)\\
                                              &\rightarrow \widetilde{H}_{n-1}\left( \left( D^n\setminus\set{0} \right)\setminus U \right)\\
                                              &\cong \widetilde{H}_{n-1}\left( S^{n-1}\sqcup S^{n-1} \right)\\
                                              &\cong \widetilde{H}_{n-1}\left( S^{n-1} \right)\oplus \widetilde{H}_{n-1}\left( S^{n-1} \right)\\
                                              &\cong \Z\oplus \Z,
  \end{align*}
  which contradicts the classification of finitely generated abelian groups.
\end{solution}
\section{Current Problems}%
\begin{problem}[Problem 1]
  For every $n\geq 1$, give an example of a surjective map $S^n\rightarrow S^n$ of degree $0$.
\end{problem}
\begin{solution}
  If we view $S^n\subseteq \R^{n+1}$, then we may consider the projection map onto the first $n$ coordinates, $\pi\colon S^n\rightarrow D^n\subseteq \R^{n+1}$, which maps $S^n$ surjectively onto a copy of $D^n$ sitting within $\R^{n+1}$.

  The induced map $\pi_{\ast}\colon H_n\left( S^n \right)\rightarrow H_n\left( D^n \right)$ is thus the zero map. Viewing $D^n$ as a CW complex structure with the boundary a copy of $S^{n-1}$ and the interior as a $n$-cell, we observe that $\left( D^n,S^{n-1} \right)$ is a CW pair, so it is good, and thus
  \begin{align*}
    H_n\left( D^n,S^{n-1} \right) &\cong H_n\left( D^{n}/S^{n-1},\ast \right)\\
                                  &\cong \widetilde{H_n}\left( S^n \right)\\
                                  &\cong H_n\left( S^n \right),
  \end{align*}
  whence upon composing with the quotient map, we have that $f\colon S^n\rightarrow S^n$ given by $q\circ\pi$ is surjective, and it induces a homomorphism from $H_n\left( S^n \right)$ to $H_n\left( S^n \right)$ that factors through the zero map, meaning that $f$ has degree $0$.
\end{solution}
\begin{problem}[Problem 2]
  Prove that if $n\geq 1$, then any map $S^n\rightarrow S^n$ is homotopic to a map with a fixed point.
\end{problem}
\begin{solution}
  Suppose $f\colon S^n\rightarrow S^n$ has no fixed points. Then, as part of the proof that $\deg\left( f \right) = \left( -1 \right)^{n+1}$, we have that $f\simeq a$, where $a$ denotes the antipodal map. Therefore, we only need to show that the antipodal map is homotopic to a map with a fixed point.

  If $n$ is odd, then the antipodal map has degree $1$, and specifically, the transformation corresponding to the antipodal map is in $ \text{SO}(n+1) $. Since $\text{SO}(n+1)$ is path-connected, there is a path from the identity to the antipodal map in $\text{SO}(n+1)$, which corresponds to a homotopy between the antipodal map and the identity map. Since the identity map preserves fixed points, we are done in this case.

  If $n$ is even, then the antipodal map has degree $-1$, and the transformation corresponding to the antipodal map is in $\text{O}\left( n+1 \right)$ but not in $\text{SO}\left( n+1 \right)$. Since the orthogonal group has two path components, it follows that the antipodal map is in the same path component as some reflection. Similarly, a path from the antipodal transformation to a reflection defines a homotopy to a reflection, and since reflections have a fixed point, we are done.
\end{solution}
\begin{problem}[Problem 3]
  Let $f_k\colon S^1\rightarrow S^1$ be the map $z\mapsto z^{k}$ for any $k\in \Z$. Prove that the local degree at each root of unity is $1$ if $k > 0$ and $-1$ if $k < 0$.
\end{problem}
\begin{solution}
  Since the map $z\mapsto z^{k}$ is a covering map of $S^1$ onto itself, it follows that for a suitably refined neighborhood of $1$ in $S^1$ (which we can take without loss of generality to be homeomorphic to an open interval), we have homeomorphic disjoint copies at each of the $\zeta_k^{j}$ where $0\leq j \leq k-1$. Locally, this gives a homeomorphism of pairs $\left( U_k,U_k\setminus\set{\zeta_k} \right)$ and $\left( V,V\setminus\set{1} \right)$. Therefore, the local degree of $f_k$ is either $1$ or $-1$, following from the commutativity of the following diagram giving the local degree of $f_k$.
  \begin{center}
    % https://tikzcd.yichuanshen.de/#N4Igdg9gJgpgziAXAbVABwnAlgFyxMJZABgBpiBdUkANwEMAbAVxiRAAkB9ARgAoBVTgGtSgoQB1xcGDgC2WMEziTpOYJIBeMusIC+AShC7S6TLnyEUZblVqMWbLnwDKAPW6k33FTPmLlUjLq4lo4OkIGRiYgGNh4BERkAEy29MysiBw8vF6GxqZxFkQelNRpDplOvABqpNU+cgpKDcDckfkxZvGWyB42ZfYZWS7unu4Nfs2Bam150bHmCSgeKQPpjtm5RrYwUADm8ESgAGYAThCySEnUOBBIAKwdZxdXN3eIbdHPlx9vSMRPc4-MggW5IADMgJeiHBf0QABYoT94XDHhRdEA
\begin{tikzcd}
{H_1(U_k,U_k\setminus\set{\zeta_k})} \arrow[r] & {H_1(V,V\setminus\set{1})} \arrow[d]     \\
{H_1(S^1,S^1\setminus\set{\zeta_k})} \arrow[u] & {H_1(S^1,S^1\setminus\set{1})} \arrow[d] \\
H_1(S^1) \arrow[r] \arrow[u]                   & H_1(S^1)                                
\end{tikzcd}
  \end{center}
  If $k > 0$, then $f_k$ is orientation-preserving about each $\zeta_k^{j}$, as it stretches the neighborhood about $\zeta_k^{j}$ by a factor of $k$ and rotates by an angle commensurate with yielding the desired neighborhood $V$ of $1$, giving that $f_k$ has degree $1$. If $k < 0$, then we compose $z\mapsto z^{-1}$; since $z^{-1} = \overline{z}$ on $S^1$, this is a reflection, so in this case, $f_k$ has degree $-1$.
\end{solution}
\end{document}
